\documentclass[12pt, a4paper]{article}

\usepackage[utf8]{inputenc}
\usepackage[T1]{fontenc}
\usepackage[brazilian]{babel}
\usepackage{geometry}
\usepackage{booktabs}
\usepackage{longtable}
\usepackage{array}
\usepackage{ragged2e}
\usepackage{xcolor}
\usepackage{titlesec}
\usepackage{xurl}
\usepackage{enumitem}
\usepackage{fancyhdr}
\usepackage{graphicx}
\usepackage{float}
\usepackage{hyperref}

\geometry{
  top=2.5cm,
  bottom=2.5cm,
  left=3cm,
  right=2.5cm
}

\setlength{\headheight}{15pt}
\setlength{\tabcolsep}{4pt}
\setlength{\emergencystretch}{2em}
\newcolumntype{L}[1]{>{\RaggedRight\arraybackslash}p{#1}}
\newcolumntype{T}[1]{>{\ttfamily\RaggedRight\arraybackslash}p{#1}}

% Cabeçalho e rodapé
\pagestyle{fancy}
\fancyhf{}
\rhead{Lendify — Módulo de Equipamentos}
\lhead{Segurança da Informação — N3}
\rfoot{Página \thepage}

% Cores
\definecolor{lendifyblue}{RGB}{30, 80, 160}
\definecolor{tableheader}{RGB}{220, 230, 245}
\definecolor{alertred}{RGB}{180, 30, 30}
\definecolor{okgreen}{RGB}{30, 130, 60}

% Estilo de títulos
\titleformat{\section}{\large\bfseries\color{lendifyblue}}{}{0em}{}[\titlerule]
\titleformat{\subsection}{\normalsize\bfseries\color{lendifyblue}}{}{0em}{}

\hypersetup{
  colorlinks=true,
  linkcolor=lendifyblue,
  urlcolor=lendifyblue,
  pageanchor=false
}

% -----------------------------------------------------------
\begin{document}

% Capa
\begin{titlepage}
  \centering
  \vspace*{2cm}
  {\Large\bfseries Centro Universitário Católica SC\\[0.3cm]}
  {\large Curso de Engenharia de Software\\[0.3cm]}
  {\large Segurança da Informação — Avaliação N3\\[2cm]}

  {\huge\bfseries\color{lendifyblue} Lendify\\[0.5cm]}
  {\Large Sistema de Reserva de Equipamentos\\[0.5cm]}
  \rule{\linewidth}{0.5pt}\\[0.5cm]
  {\LARGE\bfseries Especificação do Módulo de Equipamentos\\[0.3cm]}
  {\large Entidades, Campos, Tipos e Regras de Negócio\\[2cm]}
  \rule{\linewidth}{0.5pt}\\[1cm]

  \begin{tabular}{ll}
    \textbf{Responsável pela entrega:} & Miguel Angel Balladares Huertas \\
    \textbf{Repositório:}              & \url{https://github.com/migueldufloth/lendify} \\
    \textbf{Branch:}                   & \texttt{feature/equipamentos} \\
    \textbf{Checkpoint:}               & Fundação Técnica — 08/06/2026 \\
    \textbf{Professor:}                & Edson Vaz Lopes \\
  \end{tabular}

  \vfill
  {\large Joinville, junho de 2026}
\end{titlepage}

% Sumário
\tableofcontents
\newpage
\pagenumbering{arabic}
\setcounter{page}{1}
\hypersetup{pageanchor=true}

% -----------------------------------------------------------
\section{Introdução}

Este documento descreve a especificação técnica do módulo de \textbf{Equipamentos} do
sistema \textbf{Lendify}, desenvolvido como trabalho prático da disciplina de Segurança
da Informação (Avaliação N3) do curso de Engenharia de Software da Católica SC.

O módulo de Equipamentos é central no sistema: é o recurso que será reservado pelos
solicitantes, gerenciado pelos administradores e monitorado pelos operadores durante o
processo de empréstimo e devolução.

Este documento cobre:
\begin{itemize}
  \item a entidade principal \texttt{Equipamento} e todos os seus campos;
  \item os tipos de dados e constraints de cada campo;
  \item as regras de negócio que governam o ciclo de vida de um equipamento;
  \item as permissões de acesso por perfil de usuário;
  \item os eventos que devem ser registrados em log de auditoria.
\end{itemize}

% -----------------------------------------------------------
\section{Visão Geral do Módulo}

O módulo de Equipamentos é responsável pelo cadastro e controle de todos os itens
disponíveis para reserva no sistema, como projetores, notebooks e kits de laboratório.

Um equipamento possui um ciclo de vida baseado em \textbf{status}, que reflete seu estado
atual dentro do sistema: se está disponível para nova reserva, se já está reservado por
alguém, ou se está fora de operação por manutenção.

Apenas o perfil \textbf{Administrador} pode criar, editar ou remover equipamentos.
Os perfis \textbf{Operador} e \textbf{Solicitante} podem apenas visualizar os equipamentos
disponíveis.

% -----------------------------------------------------------
\section{Entidade: Equipamento}

\subsection{Descrição}

A entidade \texttt{Equipamento} representa um item físico cadastrado no sistema que pode
ser reservado por um solicitante. Cada equipamento possui identificação única, descrição,
categoria, número de série e um status que controla sua disponibilidade.

\subsection{Campos, Tipos e Constraints}

\begin{table}[H]
\centering
\small
\begin{tabular}{T{2.8cm} L{3.2cm} L{2.8cm} L{4.6cm}}
  \toprule
  \textbf{Campo} & \textbf{Tipo} & \textbf{Obrigatório} & \textbf{Descrição / Constraints} \\
  \midrule

  id
    & INTEGER (PK)
    & Automático
    & Chave primária. Gerado automaticamente pelo banco. Nunca editável. \\
  \addlinespace

  nome
    & VARCHAR(150)
    & Sim
    & Nome do equipamento. Mín. 3 caracteres. Ex.: \textit{Projetor Epson X41+}. \\
  \addlinespace

  descricao
    & TEXT
    & Não
    & Descrição detalhada do equipamento. Campo opcional. \\
  \addlinespace

  categoria
    & VARCHAR(100)
    & Sim
    & Categoria do equipamento. Ex.: \textit{Projetor}, \textit{Notebook}, \textit{Kit}. \\
  \addlinespace

  numero\_serie
    & VARCHAR(100)
    & Sim
    & Número de série único. Constraint \texttt{UNIQUE}. Identifica fisicamente o item. \\
  \addlinespace

  status
    & ENUM
    & Sim
    & Situação atual do equipamento. Valores permitidos: \texttt{disponivel}, \texttt{reservado}, \texttt{manutencao}. Valor padrão: \texttt{disponivel}. \\
  \addlinespace

  criado\_por
    & FK → Usuario
    & Automático
    & Chave estrangeira para o usuário (Admin) que cadastrou o equipamento. Preenchido automaticamente no momento da criação. \\
  \addlinespace

  criado\_em
    & DATETIME
    & Automático
    & Data e hora do cadastro. Preenchido automaticamente (\texttt{auto\_now\_add}). Nunca editável. \\

  \bottomrule
\end{tabular}
\caption{Campos da entidade Equipamento}
\end{table}

\subsection{Representação no Modelo Django}

O modelo será implementado no back-end Django conforme a estrutura abaixo:

\begin{verbatim}
class Equipamento(models.Model):

    STATUS_CHOICES = [
        ('disponivel',  'Disponível'),
        ('reservado',   'Reservado'),
        ('manutencao',  'Em Manutenção'),
    ]

    nome          = models.CharField(max_length=150)
    descricao     = models.TextField(blank=True, null=True)
    categoria     = models.CharField(max_length=100)
    numero_serie  = models.CharField(max_length=100, unique=True)
    status        = models.CharField(
                        max_length=20,
                        choices=STATUS_CHOICES,
                        default='disponivel'
                    )
    criado_por    = models.ForeignKey(
                        'Usuario',
                        on_delete=models.SET_NULL,
                        null=True,
                        related_name='equipamentos_criados'
                    )
    criado_em     = models.DateTimeField(auto_now_add=True)

    class Meta:
        db_table = 'equipamento'
        ordering = ['nome']

    def __str__(self):
        return f"{self.nome} [{self.numero_serie}]"
\end{verbatim}

% -----------------------------------------------------------
\section{Ciclo de Vida e Status}

O campo \texttt{status} controla a disponibilidade do equipamento para novas reservas.
As transições de status válidas são:

\begin{table}[H]
  \centering
  \begin{tabular}{ccc}
    \toprule
    \textbf{Status Atual} & \textbf{Evento que causa a transição} & \textbf{Próximo Status} \\
    \midrule
    \texttt{disponivel}  & Reserva aprovada pelo Operador/Admin  & \texttt{reservado}    \\
    \texttt{reservado}   & Devolução registrada pelo Operador/Admin & \texttt{disponivel} \\
    \texttt{disponivel}  & Admin coloca em manutenção            & \texttt{manutencao}   \\
    \texttt{manutencao}  & Admin retira de manutenção            & \texttt{disponivel}   \\
    \texttt{reservado}   & Admin força encerramento              & \texttt{disponivel}   \\
    \bottomrule
  \end{tabular}
  \caption{Transições válidas de status do Equipamento}
\end{table}

\textbf{Regra importante:} Um equipamento com status \texttt{reservado} ou
\texttt{manutencao} \textbf{não pode} receber novas reservas. O sistema deve bloquear
a criação de uma reserva se o equipamento não estiver com status \texttt{disponivel}.
Essa validação deve ocorrer \textbf{no back-end} (serializer DRF), e não apenas no front-end.

% -----------------------------------------------------------
\section{Regras de Negócio}

\subsection{RN-01 — Criação de Equipamento}
\begin{itemize}
  \item Apenas usuários com perfil \textbf{Administrador} podem criar equipamentos.
  \item Os campos \texttt{nome}, \texttt{categoria} e \texttt{numero\_serie} são obrigatórios.
  \item O \texttt{numero\_serie} deve ser único no banco de dados. Tentativa de cadastrar
        número de série duplicado deve retornar erro de validação.
  \item O campo \texttt{criado\_por} é preenchido automaticamente com o usuário autenticado
        que realizou a requisição — nunca enviado pelo cliente.
  \item O status inicial é sempre \texttt{disponivel}.
\end{itemize}

\subsection{RN-02 — Edição de Equipamento}
\begin{itemize}
  \item Apenas usuários com perfil \textbf{Administrador} podem editar equipamentos.
  \item O campo \texttt{id} e o campo \texttt{criado\_em} são imutáveis após a criação.
  \item O campo \texttt{status} pode ser editado diretamente apenas pelo Administrador
        (ex.: colocar em manutenção). Operador e Solicitante não podem alterar o status
        diretamente.
  \item A alteração de status para \texttt{manutencao} não é permitida enquanto o
        equipamento estiver com status \texttt{reservado} — deve haver devolução primeiro.
\end{itemize}

\subsection{RN-03 — Exclusão de Equipamento}
\begin{itemize}
  \item Apenas usuários com perfil \textbf{Administrador} podem excluir equipamentos.
  \item Um equipamento com status \texttt{reservado} \textbf{não pode ser excluído}
        enquanto houver reserva ativa vinculada a ele.
  \item Recomenda-se exclusão lógica (soft delete) ou manutenção do histórico de reservas
        mesmo após exclusão do equipamento.
\end{itemize}

\subsection{RN-04 — Visualização de Equipamentos}
\begin{itemize}
  \item Todos os perfis (Solicitante, Operador, Administrador) podem listar e visualizar
        equipamentos.
  \item O Solicitante visualiza apenas equipamentos com status \texttt{disponivel}
        (filtro aplicado no back-end).
  \item O Operador e o Administrador visualizam equipamentos em qualquer status.
\end{itemize}

\subsection{RN-05 — Disponibilidade para Reserva}
\begin{itemize}
  \item Somente equipamentos com status \texttt{disponivel} podem ser selecionados
        em uma nova reserva.
  \item Essa verificação é obrigatória no back-end no momento da criação da reserva.
  \item Equipamentos em \texttt{manutencao} ou \texttt{reservado} devem retornar erro
        ao tentar reservá-los.
\end{itemize}

\subsection{RN-06 — Número de Série}
\begin{itemize}
  \item O número de série identifica fisicamente o equipamento e deve ser único.
  \item Não é permitido reutilizar o número de série de um equipamento excluído.
  \item A validação de unicidade deve ser realizada no serializer (back-end),
        não apenas no front-end.
\end{itemize}

% -----------------------------------------------------------
\section{Matriz de Permissões do Módulo}

\begin{center}
\begin{tabular}{lccc}
  \toprule
  \textbf{Operação}            & \textbf{Solicitante} & \textbf{Operador} & \textbf{Admin} \\
  \midrule
  Listar equipamentos          & \textcolor{okgreen}{Sim (filtrado)} & \textcolor{okgreen}{Sim} & \textcolor{okgreen}{Sim} \\
  Ver detalhe de equipamento   & \textcolor{okgreen}{Sim}            & \textcolor{okgreen}{Sim} & \textcolor{okgreen}{Sim} \\
  Criar equipamento            & \textcolor{alertred}{Não}           & \textcolor{alertred}{Não} & \textcolor{okgreen}{Sim} \\
  Editar equipamento           & \textcolor{alertred}{Não}           & \textcolor{alertred}{Não} & \textcolor{okgreen}{Sim} \\
  Excluir equipamento          & \textcolor{alertred}{Não}           & \textcolor{alertred}{Não} & \textcolor{okgreen}{Sim} \\
  Alterar status (manutenção)  & \textcolor{alertred}{Não}           & \textcolor{alertred}{Não} & \textcolor{okgreen}{Sim} \\
  \bottomrule
\end{tabular}
\end{center}

% -----------------------------------------------------------
\section{Eventos de Log de Auditoria}

Todas as ações relevantes sobre equipamentos devem gerar um registro na tabela
\texttt{LogAuditoria}, conforme definido no modelo de dados do sistema.

\begin{table}[H]
\centering
\small
\begin{tabular}{L{3.4cm} L{3.1cm} L{6.5cm}}
  \toprule
  \textbf{Evento}                 & \textbf{Ação (campo \texttt{acao})} & \textbf{Descrição do registro} \\
  \midrule

  Cadastro de equipamento
    & \texttt{CRIAR}
    & Admin X cadastrou o equipamento Y (id, nome, número de série). \\
  \addlinespace

  Edição de equipamento
    & \texttt{EDITAR}
    & Admin X editou o equipamento Y. Campos alterados: [lista]. \\
  \addlinespace

  Exclusão de equipamento
    & \texttt{EXCLUIR}
    & Admin X excluiu o equipamento Y (id, nome, número de série). \\
  \addlinespace

  Alteração de status
    & \texttt{EDITAR}
    & Admin X alterou status do equipamento Y de \texttt{[anterior]} para \texttt{[novo]}. \\
  \addlinespace

  Tentativa de acesso negado
    & \texttt{ACESSO\_NEGADO}
    & Usuário X tentou operar o equipamento Y sem permissão. \\

  \bottomrule
\end{tabular}
\caption{Eventos de log de auditoria do módulo de Equipamentos}
\end{table}

% -----------------------------------------------------------
\section{Validações no Back-end (Serializer DRF)}

As seguintes validações devem ser implementadas no serializer Django REST Framework
do módulo de Equipamentos:

\begin{itemize}
  \item \textbf{nome}: mínimo de 3 caracteres, máximo de 150.
  \item \textbf{categoria}: campo obrigatório, não pode ser vazio.
  \item \textbf{numero\_serie}: obrigatório, único, máximo de 100 caracteres. Validar
        unicidade via \texttt{UniqueValidator}.
  \item \textbf{status}: deve ser um dos valores permitidos: disponível, reservado
        ou manutenção.
  \item \textbf{criado\_por}: campo somente leitura — nunca aceito na entrada da
        requisição. Preenchido automaticamente via \texttt{perform\_create}.
  \item \textbf{criado\_em}: campo somente leitura — nunca aceito na entrada da
        requisição.
\end{itemize}

% -----------------------------------------------------------
\section{Dados Fictícios para Demonstração}

Para demonstração e testes do sistema, os seguintes equipamentos fictícios serão
cadastrados via \textit{fixture} ou script de seed:

\begin{center}
\begin{tabular}{llll}
  \toprule
  \textbf{Nome}              & \textbf{Categoria} & \textbf{Nº de Série} & \textbf{Status} \\
  \midrule
  Projetor Epson X41+        & Projetor           & SN-PRJ-001           & disponivel \\
  Notebook Dell Latitude     & Notebook           & SN-NTB-001           & disponivel \\
  Notebook Dell Latitude     & Notebook           & SN-NTB-002           & reservado  \\
  Kit Arduino Iniciante      & Kit                & SN-KIT-001           & disponivel \\
  Projetor BenQ MH560        & Projetor           & SN-PRJ-002           & manutencao \\
  Notebook Lenovo ThinkPad   & Notebook           & SN-NTB-003           & disponivel \\
  Kit Robótica Intermediário & Kit                & SN-KIT-002           & disponivel \\
  \bottomrule
\end{tabular}
\end{center}

Todos os dados acima são \textbf{estritamente fictícios}, conforme exigido pelas
regras da avaliação.

% -----------------------------------------------------------
\section{Considerações de Segurança}

\begin{itemize}
  \item \textbf{Autorização por perfil:} todas as rotas de criação, edição e exclusão
        de equipamentos devem verificar se o usuário autenticado possui perfil
        \texttt{admin}. Caso contrário, retornar HTTP 403 e registrar o evento em log.

  \item \textbf{Validação no back-end:} nenhuma validação pode depender exclusivamente
        do front-end. O serializer deve rejeitar dados inválidos independentemente
        da interface.

  \item \textbf{Campo criado\_por:} nunca confiar no valor enviado pelo cliente.
        O campo deve ser preenchido a partir do token JWT do usuário autenticado.

  \item \textbf{Número de série:} dado sensível de rastreabilidade do patrimônio.
        Não deve ser exposto em endpoints públicos ou sem autenticação.

  \item \textbf{Log de acesso negado:} qualquer tentativa de operação não autorizada
        sobre equipamentos deve gerar registro no log de auditoria, incluindo IP
        e timestamp.
\end{itemize}

% -----------------------------------------------------------
\section{Referências}

\begin{itemize}
  \item Repositório do projeto: \url{https://github.com/migueldufloth/lendify}
  \item Documento da Avaliação N3 — Prof. Edson Vaz Lopes — Católica SC, junho de 2026.
  \item Django REST Framework — Serializers: \url{https://www.django-rest-framework.org/api-guide/serializers/}
  \item Django Models: \url{https://docs.djangoproject.com/en/stable/topics/db/models/}
\end{itemize}

\end{document}